\section{Reference protocol and method definitions}
\label{app:protocol}
The comparisons that support our final selection use the protocol below. The three retained methods and the plain control are defined as they were executed in the final batch. The same definitions are frozen in our shared code, which reproduces the executed intervention states and parameter updates bitwise on CPU at sampled updates.

\subsection{Training recipe}
\begin{center}\small
\begin{tabularx}{\linewidth}{@{}lX@{}}
\toprule
Data & CIFAR-10, all 50,000 training and 10,000 test images; no augmentation; per-channel normalization with the statistics of the training images.\\
Model & ResNet-20 with BatchNorm, widths 16/32/64, option-A shortcuts (subsampling and zero-padding), 269,722 parameters; pinned initial weights per seed.\\
Optimizer & SGD, learning rate $\eta=0.005$, momentum $0.9$, weight decay $5\times10^{-4}$ on all parameters; $\eta\,(u+1)/60$ for the first 60 updates, then cosine decay to zero over the global update index $u$.\\
Batches & A pinned permutation per epoch; batches of 128 images (the last one of 80), each split into microbatches of 32 whose losses are weighted by their number of examples. BatchNorm therefore normalizes over microbatches.\\
Budget & 30 epochs of 391 updates, 11,730 updates in total.\\
Evaluation & Before training and after every epoch, on the test set and on a fixed probe of 500 training images, in eval mode with BatchNorm running statistics.\\
\bottomrule
\end{tabularx}
\end{center}

\subsection{Conventions}
\label{app:conventions}
\paragraph{Statistics and pairing.} We report the record after epoch 30, without best-epoch selection, as the mean and sample standard deviation (SD) over seeds; a configuration with a single valid seed has no SD. Within a batch, runs with the same seed share their initial weights, BatchNorm buffers and data order, and we compute gains as means of per-seed paired differences. We do not pair runs across batches: plain runs on identical pinned assets finished up to $\plainSpreadMax$ points apart in different batches, for reasons we have not established. Bars and bands in the figures show $\pm1$ sample SD; they are not confidence intervals.

\paragraph{Evaluation paths.} The \emph{current path} evaluates the network with the intervention of the epoch just completed. The \emph{target path} evaluates the same weights at the native resolution without added filters, which is the network kept at the end. After epoch 30 the two paths coincide for every retained method. Before that, the target path uses BatchNorm running statistics accumulated on the current path, so it combines the effect of removing the intervention with a mismatch of these statistics.

\paragraph{Training loss.} Training cross-entropy is evaluated on the fixed probe of 500 training images, not on the full training set, which we did not evaluate. We also record the mean training loss of each epoch, computed while the weights are being updated.

\paragraph{Records.} Run identifiers, asset hashes, commits and the complete configuration catalogue, including failed runs, are archived with the sources of this report (\texttt{paper/archive/}).

\subsection{Insertion sites and the reduction point}
\label{app:sites}
``After block 1'' denotes the complete output of the second residual block, after its residual addition and final ReLU. The adaptive max-pooling is applied to this tensor as it enters the third block, so both branches of every later block, including the shortcut, see the reduced grid. Table~\ref{tab:sites} lists the sites in forward order.

\begin{table}[!htbp]
\caption{Insertion sites of the ResNet-20 reference model. Each residual block has two convolutions, and stride-2 convolutions open blocks 3 and 6. Map sides are given for $r=16/24/32$ at the reduction point; the Gaussian-only method never reduces, so its maps are those of the $r=32$ column. Post-ReLU sites are the stem ReLU output ($p_0$) and the output of each block ($p_1$--$p_9$).}
\label{tab:sites}
\centering\small
\begin{tabular}{@{}lccc@{}}
\toprule
Stage & Map side, $r=16/24/32$ & Conv outputs (\methodRG) & Post-ReLU (\methodG)\\
\midrule
Stem & 32 / 32 / 32 & $c_0$ & $p_0$\\
Blocks 0--1 & 32 / 32 / 32 & $c_1$--$c_4$ & $p_1,p_2$\\
\emph{Reduction (input of block 2)} & 16 / 24 / 32 & --- & ---\\
Block 2 & 16 / 24 / 32 & $c_5,c_6$ & $p_3$\\
Blocks 3--5 & 8 / 12 / 16 & $c_7$--$c_{12}$ & $p_4$--$p_6$\\
Blocks 6--8 & 4 / 6 / 8 & $c_{13}$--$c_{18}$ & $p_7$--$p_9$\\
\bottomrule
\end{tabular}
\end{table}

In the combined method, each Gaussian sits between a convolution and its BatchNorm, so BatchNorm normalizes filtered activations. In the Gaussian-only method, the filtered tensor is the block output that the next block, or the pooling layer for $p_9$, receives on both paths; the ReLU inside each block is not a site.

\subsection{Schedules, operators and bypass}
\label{app:schedules}
% GENERATED by paper/tools/make_paper_assets.py -- do not edit by hand.
% Sources: continuation_core.methods (PLATEAU_G, RPROG, METHODS) and continuation_core.controller.InterventionController.per_site_sigma
% Provenance (configs, seeds, cells, files): paper/provenance/generated_assets.json
\begin{table}[!htbp]
\caption{Executed schedules, evaluated through the controller of the frozen presets. Epochs are zero-based; a state holds for a whole epoch of 391 updates. $r$ is the side of the map entering residual block index 2 (\methodR{} and \methodRG), $G$ the reference Gaussian level. Effective $\sigma$ is in pixels of the feature map at the site and is the same at every site of a method.}
\label{tab:schedules}
\centering\small
\begin{tabular}{@{}rrcccc@{}}
\toprule
Epochs & First update & $r$ & $G$ & \methodG: $\sigma$ (10 sites) & \methodRG: $\sigma=\frac{r}{32}G$ (19 sites)\\
\midrule
0--2 & 0 & 16 & 1.00 & 1.000 & 0.500\\
3--5 & 1173 & 16 & 0.85 & 0.850 & 0.425\\
6--8 & 2346 & 24 & 0.70 & 0.700 & 0.525\\
9--11 & 3519 & 24 & 0.60 & 0.600 & 0.450\\
12--14 & 4692 & 32 (bypass) & 0.50 & 0.500 & 0.500\\
15--17 & 5865 & 32 (bypass) & 0.40 & 0.400 & 0.400\\
18--20 & 7038 & 32 (bypass) & 0.30 & 0.300 & 0.300\\
21--29 & 8211 & 32 (bypass) & 0.00 & 0 (bypass) & 0 (bypass)\\
\bottomrule
\end{tabular}
\end{table}


A state holds for a whole epoch: a transition at epoch $e$ means that update $391e$ is the first to use the new state, and the record written after $k$ completed epochs evaluates the current path under the state of epoch $k-1$.

The Gaussian kernel has the same nine taps for every $\sigma$, $w_d\propto\exp(-d^2/(2\sigma^2))$ for $d=-4,\dots,4$, normalized to sum one. It is applied separably and channel by channel, with reflection padding that does not repeat the edge sample; on $4\times4$ maps, where the radius equals the map side, the same reflection is applied by explicit indexing. The scale $\sigma$ is measured in pixels of the feature map at the site.

The combined method uses $\sigma=(r/32)\,G$ at all 19 sites, including $c_0$--$c_4$, which remain at $32\times32$ upstream of the reduction. This rescaling comes from our earlier input-resolution arms (Appendix~\ref{app:explore-resolution}) and makes the effective scale non-monotone: it increases from $0.425$ to $0.525$ when $r$ changes from 16 to 24 at epoch 6. We keep this executed rule rather than a local width ratio, so the combined method is not the literal composition of the two other methods.

When $G=0$, a Gaussian hook returns its input unchanged, and when $r=32$ the reduction hook returns its input. In the final state $(r=32,G=0)$, every method therefore computes exactly the plain ResNet-20. Neither BatchNorm statistics nor optimizer state are reset at a transition.
